
\documentclass{article}
\usepackage{colortbl}
\usepackage{makecell}
\usepackage{multirow}
\usepackage{supertabular}

\begin{document}

\newcounter{utterance}

\centering \large Interaction Transcript for game `cladder', experiment `full\_v1.5\_default', episode 905 with score{-}Qwen3.5{-}9B{-}sp.
\vspace{24pt}

{ \footnotesize  \setcounter{utterance}{1}
\setlength{\tabcolsep}{0pt}
\begin{supertabular}{c@{$\;$}|p{.15\linewidth}@{}p{.15\linewidth}p{.15\linewidth}p{.15\linewidth}p{.15\linewidth}p{.15\linewidth}}
    \# & \multicolumn{2}{c}{Player} && \multicolumn{2}{c}{Game Master} \\
    \hline

    \theutterance \stepcounter{utterance}  
    & & & \multicolumn{4}{p{0.6\linewidth}}{
        \cellcolor[rgb]{0.9,0.9,0.9}{
            \makecell[{{p{\linewidth}}}]{
                \texttt{\tiny{[P1$\langle$GM]}}
                \texttt{You are an expert in causal inference. The following question is not a typical commonsense query, but rather a meticulously designed question created by a professor specializing in causal inference, intended to assess the students' mastery of the course content.} \\
\texttt{Imagine a self{-}contained, hypothetical world with only the following conditions, and without any unmentioned factors or causal relationships: Smoking has a direct effect on effort and college admission. Effort has a direct effect on college admission. For nonsmokers, the probability of college admission is 76\%. For smokers, the probability of college admission is 22\%. For smokers, would it be less likely to see college admission if the student had been a nonsmoker?} \\
\texttt{Start your answer with 'Yes' or 'No', followed by additional reasoning or evidence to support your explanation.} \\
            }
        }
    }
     \\ \\

    \theutterance \stepcounter{utterance}  
    & \multicolumn{4}{p{0.6\linewidth}}{
        \cellcolor[rgb]{0.561,0.737,0.561}{
            \makecell[{{p{\linewidth}}}]{
                \texttt{\tiny{[P1$\rangle$GM]}}
                \texttt{No} \\
\\ 
\texttt{**Reasoning:**} \\
\\ 
\texttt{To determine the answer, we must distinguish between the **observed association** (what we see in the data) and the **causal effect** (what would happen if we intervened).} \\
\\ 
\texttt{1.  **The Observed Data:** The prompt states that for smokers, the probability of college admission is 22\%, while for nonsmokers, it is 76\%. If we simply compare these two numbers, it appears that being a nonsmoker is associated with a much higher likelihood of admission.} \\
\\ 
\texttt{2.  **The Causal Structure:** The problem defines a specific causal graph:} \\
\texttt{    *   Smoking \$\textbackslash{}rightarrow\$ Effort} \\
\texttt{    *   Smoking \$\textbackslash{}rightarrow\$ College Admission} \\
\texttt{    *   Effort \$\textbackslash{}rightarrow\$ College Admission} \\
\\ 
\texttt{    Crucially, the prompt states there are "no unmentioned factors." This implies there is no **confounding variable** (such as socioeconomic status or innate ability) that causes both smoking and low effort/admission. In a confounded scenario, smokers might smoke *because* they have low effort or low ability, which also lowers their admission chances. However, in this hypothetical world, the only link between smoking and admission is the direct effect of smoking itself and the indirect effect via effort.} \\
\\ 
\texttt{3.  **The Counterfactual Question:** The question asks: "For smokers, would it be less likely to see college admission if the student had been a nonsmoker?"} \\
\texttt{    *} \\
            }
        }
    }
    & & \\ \\

    \theutterance \stepcounter{utterance}  
    & & & \multicolumn{2}{p{0.3\linewidth}}{
        \cellcolor[rgb]{0.9,0.9,0.9}{
            \makecell[{{p{\linewidth}}}]{
                \texttt{\tiny{[GM$|$GM]}}
                \texttt{No} \\
            }
        }
    }
    & & \\ \\

    \theutterance \stepcounter{utterance}  
    & & & \multicolumn{2}{p{0.3\linewidth}}{
        \cellcolor[rgb]{0.9,0.9,0.9}{
            \makecell[{{p{\linewidth}}}]{
                \texttt{\tiny{[GM$|$GM]}}
                \texttt{no} \\
            }
        }
    }
    & & \\ \\

    \theutterance \stepcounter{utterance}  
    & & & \multicolumn{2}{p{0.3\linewidth}}{
        \cellcolor[rgb]{0.9,0.9,0.9}{
            \makecell[{{p{\linewidth}}}]{
                \texttt{\tiny{[GM$|$GM]}}
                \texttt{game\_result = WIN} \\
            }
        }
    }
    & & \\ \\

\end{supertabular}
}

\end{document}
